\documentclass[conference]{IEEEtran}
\IEEEoverridecommandlockouts
\usepackage{cite}
\usepackage{amsmath,amssymb,amsfonts}
\usepackage{algorithmic}
\usepackage{graphicx}
\usepackage{textcomp}
\usepackage{xcolor}
\usepackage{booktabs}
\def\BibTeX{{\rm B\kern-.05em{\sc i\kern-.025em b}\kern-.08em
    T\kern-.1667em\lower.7ex\hbox{E}\kern-.125emX}}
\begin{document}

\title{ScholarPilot: A Mobile Platform for Embedding-Based Paper Recommendation, Abstract-Grounded AI Assistance and Location-Aware Conference Discovery}

\author{
    \centering
    \begin{tabular}{@{}c@{\hskip 1.5cm}c@{}}
        \begin{tabular}{c}
             Atharva Jadhav \\
            \textit{Department of Computer Engineering} \\
            \textit{Fr. Conceicao Rodrigues Institute of Technology}\\
            Navi Mumbai, India \\
            jadhavatharva907@gmail.com
        \end{tabular}
        &
        \begin{tabular}{c}
            Rugved Hinge \\
            \textit{Department of Computer Engineering} \\
            \textit{Fr. Conceicao Rodrigues Institute of Technology}\\
            Navi Mumbai, India \\
            shreehinge007@gmail.com
        \end{tabular}
    \end{tabular}\\[1.5em]
    \begin{tabular}{@{}c@{\hskip 1.5cm}c@{}}
        \begin{tabular}{c}
            Dr. Rakhi Kalantri \\
            \textit{Department of Computer Engineering} \\
            \textit{Fr. Conceicao Rodrigues Institute of Technology}\\
            Navi Mumbai, India\\
            rakhi.kalantri@fcrit.ac.in
        \end{tabular}
        &
        \begin{tabular}{c}
            Mrs. Smita Rukhande \\
            \textit{Department of Computer Engineering} \\
            \textit{Fr. Conceicao Rodrigues Institute of Technology}\\
            Navi Mumbai, India\\
            smita.rukhande@fcrit.ac.in
        \end{tabular}
    \end{tabular}\\[1.5em]
    \begin{tabular}{c}
        Dr. Pravin Rahate \\
        \textit{Department of Computer Engineering} \\
        \textit{Fr. Conceicao Rodrigues Institute of Technology}\\
        Navi Mumbai, India\\
        pravin.rahate@fcrit.ac.in
    \end{tabular}
}

\maketitle

\begin{abstract}
Researchers usually discover papers, understand them and find relevant events in three separate tools. This paper presents \textbf{ScholarPilot}, an Android application that brings these steps together. A content-based recommender ranks papers by the cosine similarity between SPECTER2 document embeddings and a user profile built from selected interests and bookmarks, and then diversifies the batch with Maximal Marginal Relevance. A conversational assistant answers questions about a paper using a large language model grounded in the paper's metadata and abstract. A location-aware module combines open conference listings with on-device geocoding, so that upcoming conferences are sorted by their distance from the user, shown on an OpenStreetMap map and linked to turn-by-turn directions in Google Maps. The interface follows a soft, tactile claymorphic style. We describe the architecture, the data sources and the algorithms, and we define an evaluation protocol based on Precision@K, MAP, latency and the System Usability Scale.
\end{abstract}

\begin{IEEEkeywords}
paper recommendation, document embeddings, retrieval-augmented generation, conference discovery, geolocation, mobile application
\end{IEEEkeywords}

\section{Introduction}
\subsection{Motivation}
Scholarly output keeps growing, and a student beginning a literature review has to locate, filter and read far more material than can be handled by hand. General search engines return ranked lists in response to keywords and leave the filtering effort to the reader \cite{beel2016}. Reference managers store citations but rarely help with understanding a paper, and tools that list conferences are separate from both. As a result, discovering a paper, understanding it, and finding a venue to present related work are three disconnected activities.

\subsection{Contributions}
ScholarPilot addresses this gap within a single mobile application:
\begin{itemize}
    \item a recommender that ranks candidate papers by cosine similarity between SPECTER2 embeddings \cite{cohan2020,singh2023} and a user profile, then diversifies the result with MMR \cite{carbonell1998};
    \item an assistant that answers questions about a paper with a large language model \cite{gemini2023} grounded in the paper's abstract, following the retrieval-augmented principle \cite{lewis2020};
    \item a conference module built entirely from live open data, which geocodes venues on the device, computes great-circle distances \cite{sinnott1984} and hands the user over to Google Maps for navigation;
    \item a claymorphic interface designed for visual comfort \cite{norman2013} and evaluated with the System Usability Scale \cite{brooke1996,bangor2008}.
\end{itemize}

\subsection{Paper Organization}
Section~II surveys related work. Section~III describes the architecture, Section~IV the methods, and Section~V the implementation. Section~VI defines the evaluation metrics, Section~VII reports observations, and Section~VIII concludes with future work.

\section{Literature Survey}
\subsection{Paper Recommendation}
Beel et~al.\ surveyed research-paper recommender systems and found content-based filtering to be the most widely applied approach, while evaluation practice and reproducibility remained weak \cite{beel2016}. Content-based methods describe each item by features and match them against a user profile \cite{lops2011}; collaborative and hybrid techniques are collected in the Recommender Systems Handbook \cite{ricci2015}. Classical bag-of-words weighting is a common baseline but cannot capture similarity between differently worded texts \cite{manning2008}.

\subsection{Neural Document Representations}
Transformers \cite{vaswani2017} and pre-trained encoders such as BERT \cite{devlin2019} changed how text is represented. Sentence-BERT made it practical to compare whole sentences with cosine similarity \cite{reimers2019}. For scientific papers, SPECTER learns document-level vectors from the citation graph \cite{cohan2020}, and SPECTER2 extends this and is evaluated on the SciRepEval benchmark \cite{singh2023}. Semantic Scholar serves these embeddings for its corpus through an open API \cite{kinney2023}, which is the source ScholarPilot uses. Dense representations have also been shown to outperform sparse ones for retrieval in question answering \cite{karpukhin2020}.

\subsection{Diversification}
A ranking by similarity alone tends to return near-duplicates of the strongest interest. Maximal Marginal Relevance trades relevance against redundancy by repeatedly selecting the item with the best score that is least similar to those already chosen \cite{carbonell1998}.

\subsection{LLM Summarization and Grounding}
Large language models are now the dominant approach to abstractive summarization \cite{zhang2024}, but they can state things the source does not support. Retrieval-augmented generation reduces this by conditioning the answer on supplied text \cite{lewis2020,gao2023}. ScholarPilot applies the same idea at the scale of one paper: the model is instructed to answer only from the supplied metadata and abstract and to say so when something is not there.

\subsection{Usability}
Norman's principles of affordance and feedback \cite{norman2013} motivate soft, elevated and clearly pressable controls. The System Usability Scale provides a short and widely used measure of perceived usability \cite{brooke1996}, and its scores can be interpreted on an adjective scale \cite{bangor2008}.

\subsection{Identified Gap}
Recommendation, summarization and event discovery are normally offered as separate products, and conference listings are rarely connected to where the user actually is. ScholarPilot combines the three in one mobile workflow.

\section{System Architecture}
ScholarPilot is a native Android application organised in four layers, shown in Fig.~\ref{fig:arch}.

\begin{enumerate}
    \item \textbf{Presentation layer.} A Jetpack Compose interface with onboarding, interest selection, the paper feed, search, a conference list and map, and a per-paper assistant. All components follow a claymorphic style: soft elevation, layered shadows and rounded containers.
    \item \textbf{Application layer.} A view model and repository coordinate the screens, Firebase Authentication and Firestore (accounts, bookmarks, collections and interests) and a local Room database (assistant chat history).
    \item \textbf{Intelligence layer.} A recommendation engine ranks papers using SPECTER2 vectors, and an assistant client sends a grounded prompt to the Gemini API.
    \item \textbf{Data and geolocation layer.} Retrofit clients for the Semantic Scholar and OpenAlex APIs; parsers for open conference listings; a geocoder; the device location service; the OpenStreetMap Overpass API for libraries; and Google Maps for navigation.
\end{enumerate}

% ------------------------------------------------------------------
% FIGURE 1: SYSTEM ARCHITECTURE DIAGRAM  (replace the box with \includegraphics)
% ------------------------------------------------------------------
\begin{figure*}[!t]
    \centering
    % \includegraphics[width=0.9\textwidth]{system_architecture.png}
    \fbox{\parbox[c][7.5cm][c]{0.85\textwidth}{\centering
        \textbf{[Insert System Architecture Diagram here]}\\[0.5em]
        \textit{file: system\_architecture.png}}}
    \caption{Layered architecture of ScholarPilot and the external services it uses.}
    \label{fig:arch}
\end{figure*}

\section{Methodology}
\subsection{Paper Retrieval and Candidate Generation}
Papers are obtained from the Semantic Scholar Academic Graph API \cite{kinney2023}, with the keyless OpenAlex API as a fallback when the primary service is unavailable. Each selected interest is issued as a search query, and the top results are kept together with their SPECTER2 vectors. When the user has bookmarks, the Semantic Scholar recommendation endpoint is also queried with those papers as positive seeds, and the returned papers are added to the candidate pool.

\subsection{Profile Construction and Ranking}
Let $\mathbf{b}$ be the normalised centroid of the embeddings of the user's most recent bookmarks, and let $\mathbf{c}_j$ be the normalised centroid of the embeddings returned for interest $j$. The profile is
\begin{equation}
    \mathbf{u} = \frac{\lambda_b\,\mathbf{b} + (1-\lambda_b)\,\bar{\mathbf{c}}}{\lVert \lambda_b\,\mathbf{b} + (1-\lambda_b)\,\bar{\mathbf{c}} \rVert},
\end{equation}
where $\bar{\mathbf{c}}$ is the normalised mean of the interest centroids and $\lambda_b=0.6$. When there are no bookmarks the profile is simply $\bar{\mathbf{c}}$ (cold start). Each candidate $i$ with embedding $\mathbf{d}_i$ receives the score
\begin{equation}
    s_i = \frac{\mathbf{u}\cdot\mathbf{d}_i}{\lVert\mathbf{u}\rVert\,\lVert\mathbf{d}_i\rVert} + \gamma\,\frac{\ln(1+n_i)}{\ln(1+n_{\max})},
\end{equation}
where $n_i$ is the citation count and $\gamma=0.04$ is a small prior that breaks ties in favour of established work. Papers the user has already bookmarked are removed. Candidates without an embedding fall back to a lexical overlap score between their title and the interest labels.

\subsection{Diversification}
The final batch of $K=30$ papers is chosen greedily with Maximal Marginal Relevance \cite{carbonell1998}:
\begin{equation}
    \mathrm{MMR} = \arg\max_{i \in R\setminus S}\Big[\lambda\, s_i - (1-\lambda)\max_{j\in S}\cos(\mathbf{d}_i,\mathbf{d}_j)\Big],
\end{equation}
with $\lambda=0.75$, where $S$ is the set already selected and $R$ the candidate set.

\subsection{Two-Phase Feed}
Embedding retrieval needs several rate-limited API calls. To keep the screen responsive, the feed is produced in two phases: a single keyword search fills the list immediately, and the embedding-ranked batch replaces it as soon as it is ready.

\subsection{Abstract-Grounded Assistant}
For each paper the assistant sends the title, authors, venue, field and abstract in a system instruction, together with the last ten chat messages. The instruction asks the model to answer concisely from the abstract and to say so when a detail is not present. This applies the retrieval-augmented principle \cite{lewis2020} at the level of one document.

\subsection{Conference Data and Geocoding}
Conference data is assembled from three open sources (Table~\ref{tab:sources}). Only events that start in the future are kept, and the time left to the next submission deadline is computed on the device. Venue names are converted to coordinates by the Open-Meteo geocoding service with OpenStreetMap Nominatim as a fallback, and each result is cached on the device so that every place is resolved only once. The distance from the user to a venue is the great-circle distance given by the haversine formula \cite{sinnott1984},
\begin{equation}
    d = 2R\arcsin\sqrt{\sin^2\!\frac{\Delta\varphi}{2}+\cos\varphi_1\cos\varphi_2\sin^2\!\frac{\Delta\lambda}{2}},
\end{equation}
where $R$ is the Earth's radius and $\varphi,\lambda$ are latitude and longitude. Venues in the user's own country are geocoded first, so that nearby results appear before distant ones.

\begin{table}[!t]
\caption{Open Conference Data Sources}
\label{tab:sources}
\centering
\small
\begin{tabular}{@{}p{2.4cm}p{2.9cm}p{2.1cm}@{}}
\toprule
\textbf{Source} & \textbf{Coverage} & \textbf{Fields} \\
\midrule
CCF Deadlines & Ranked computer-science conferences & dates, place, deadlines, rank \\
WikiCFP & Calls for papers by topic (11 topics) & dates, place, deadline \\
conferenceindex.org & Per-country listings of the user's country & date, city \\
\bottomrule
\end{tabular}
\end{table}

\subsection{Maps and Navigation}
The conference map is rendered with osmdroid on OpenStreetMap tiles, with markers for conferences, nearby libraries (retrieved from the Overpass API) and the user. Each conference card contains a directions button that opens Google Maps with the venue as the destination, so that routing and turn-by-turn navigation are delegated to a dedicated service.

% ------------------------------------------------------------------
% FIGURE 2: WORKFLOW / DATA-FLOW DIAGRAM  (replace the box with \includegraphics)
% ------------------------------------------------------------------
\begin{figure}[!t]
    \centering
    % \includegraphics[width=\columnwidth]{workflow_diagram.png}
    \fbox{\parbox[c][9cm][c]{0.9\columnwidth}{\centering
        \textbf{[Insert Workflow Diagram here]}\\[0.5em]
        \textit{file: workflow\_diagram.png}}}
    \caption{End-to-end workflow: onboarding, recommendation, assistant and nearby-conference discovery.}
    \label{fig:workflow}
\end{figure}

\section{Implementation}
The application is written in Kotlin with Jetpack Compose. Retrofit and OkHttp provide the network layer, with a global throttle that keeps Semantic Scholar requests within the one-request-per-second limit of an API key, and a retry on rate-limit responses. Secrets such as API keys are read from a local, version-control-ignored properties file at build time. User accounts, bookmarks, collections and interests are stored in Firebase Authentication and Cloud Firestore, while assistant conversations are kept in a local Room database. Conference and geocoding results are cached on the device for 24 hours and indefinitely respectively, which makes repeat visits instantaneous and keeps the application polite towards the free services it depends on. Location is obtained through the fused location provider after the user grants permission. Fig.~\ref{fig:workflow} shows how a request flows through these components.

% ------------------------------------------------------------------
% FIGURE 3: UI SCREENSHOTS  (replace the box with \includegraphics)
% ------------------------------------------------------------------
\begin{figure}[!t]
    \centering
    % \includegraphics[width=\columnwidth]{ui_screenshots.png}
    \fbox{\parbox[c][5.5cm][c]{0.9\columnwidth}{\centering
        \textbf{[Insert UI Screenshots here]}\\[0.5em]
        \textit{file: ui\_screenshots.png}\\
        (Home feed, assistant, conference list, conference map)}}
    \caption{Screens of ScholarPilot.}
    \label{fig:ui}
\end{figure}

\section{Evaluation Metrics}
\begin{enumerate}
    \item \textbf{Recommendation quality.} Precision@$K$ with $K=10$ and Mean Average Precision (MAP), computed over a set of test profiles with relevance judged by the participants \cite{manning2008}. The ranked feed is compared against the keyword baseline that is shown first.
    \item \textbf{Diversity.} The mean pairwise cosine similarity inside a batch, with and without MMR.
    \item \textbf{Latency.} Time to first feed render, time to the ranked feed, assistant response time, and time from granting location to the first nearby result.
    \item \textbf{Geocoding coverage.} The fraction of venue strings that resolve to coordinates.
    \item \textbf{Usability.} The System Usability Scale score \cite{brooke1996}, interpreted with established benchmarks \cite{bangor2008}, and a short questionnaire on visual comfort.
\end{enumerate}

\section{Results and Discussion}
\subsection{Data Coverage}
Table~\ref{tab:coverage} reports the amount of live data available when the system was exercised on 11~October 2026. These figures show the scale of the corpus the location module works on.

\begin{table}[!t]
\caption{Live Data Observed on 11 October 2026}
\label{tab:coverage}
\centering
\small
\begin{tabular}{@{}lc@{}}
\toprule
\textbf{Item} & \textbf{Value} \\
\midrule
Upcoming ranked CS conferences (CCF Deadlines) & 275 \\
Upcoming events listed for India (conferenceindex.org) & 460 \\
Distinct venue cities among those events & 17 \\
Embedding dimension (SPECTER2) & 768 \\
\bottomrule
\end{tabular}
\end{table}

\subsection{Recommendation and Responsiveness}
On first use the keyword feed was displayed after roughly twelve seconds and was replaced by the ranked feed after roughly forty seconds in a single informal run on an emulator; repeat visits are faster because interest vectors and conference data are cached. Feeds generated for a profile of neuroscience, chemistry and climate science were dominated by papers on those topics. A controlled measurement of Precision@10 and MAP with participants is still required and is reported in Table~\ref{tab:results}.

\begin{table}[!t]
\caption{Evaluation Summary (to be completed with measured values)}
\label{tab:results}
\centering
\begin{tabular}{@{}lcc@{}}
\toprule
\textbf{Metric} & \textbf{Value} & \textbf{Target} \\
\midrule
Precision@10         & -- & $\geq 0.70$ \\
MAP                  & -- & $\geq 0.60$ \\
Intra-batch similarity (MMR) & -- & lower is better \\
Assistant latency (s) & -- & $\leq 5$ \\
First nearby result (s) & -- & $\leq 10$ \\
SUS score            & -- & $\geq 68$ \\
\bottomrule
\end{tabular}
\end{table}

\subsection{Limitations}
Recommendation quality depends on the embeddings served for each paper, and a paper without an embedding is ranked by a weaker lexical score. The assistant sees only the abstract and not the full text, so questions about methods or results beyond the abstract cannot be answered. Conference listings are community maintained and mix prestigious and commercial events; only the CCF source carries a quality rank. Geocoding depends on how precisely a venue is written, and a venue that cannot be resolved is not shown on the map.

\section{Conclusion}
ScholarPilot places paper recommendation, grounded assistance and location-aware conference discovery in one mobile application. The recommender applies document embeddings, a profile built from interests and bookmarks, and diversified ranking, and the conference module is built from live open data with on-device geocoding, distance sorting, a map and navigation hand-off. The modular design allows each component to be improved independently.

\section{Future Scope}
\begin{itemize}
    \item \textbf{Full-text understanding:} chunked retrieval over the complete paper so that the assistant can answer questions on methods and results, and multi-modal handling of figures and tables.
    \item \textbf{Learning from behaviour:} incorporating reading time and explicit feedback into the profile, and on-device or federated updates to preserve privacy.
    \item \textbf{Richer event data:} deadline information for all sources, ranking of non-CS conferences, and notifications as submission deadlines approach.
    \item \textbf{Collaboration and offline use:} shared reading lists and annotations, and cached papers that can be queried without a connection.
\end{itemize}

\begin{thebibliography}{99}
\bibitem{beel2016}
J. Beel, B. Gipp, S. Langer, and C. Breitinger, ``Research-paper recommender systems: A literature survey,'' \textit{Int. J. Digit. Libr.}, vol. 17, no. 4, pp. 305--338, 2016.

\bibitem{gao2023}
Y. Gao \textit{et al.}, ``Retrieval-augmented generation for large language models: A survey,'' \textit{arXiv preprint arXiv:2312.10997}, 2023.

\bibitem{zhang2024}
H. Zhang, P. S. Yu, and J. Zhang, ``A systematic survey of text summarization: From statistical methods to large language models,'' \textit{arXiv preprint arXiv:2406.11289}, 2024.

\bibitem{reimers2019}
N. Reimers and I. Gurevych, ``Sentence-BERT: Sentence embeddings using Siamese BERT-networks,'' in \textit{Proc. Conf. Empirical Methods Natural Lang. Process. (EMNLP)}, 2019, pp. 3982--3992.

\bibitem{manning2008}
C. D. Manning, P. Raghavan, and H. Sch\"utze, \textit{Introduction to Information Retrieval}. Cambridge, U.K.: Cambridge Univ. Press, 2008.

\bibitem{norman2013}
D. Norman, \textit{The Design of Everyday Things: Revised and Expanded Edition}. New York, NY, USA: Basic Books, 2013.

\bibitem{lewis2020}
P. Lewis \textit{et al.}, ``Retrieval-augmented generation for knowledge-intensive NLP tasks,'' in \textit{Advances in Neural Information Processing Systems (NeurIPS)}, vol. 33, 2020, pp. 9459--9474.

\bibitem{vaswani2017}
A. Vaswani \textit{et al.}, ``Attention is all you need,'' in \textit{Advances in Neural Information Processing Systems (NeurIPS)}, vol. 30, 2017, pp. 5998--6008.

\bibitem{devlin2019}
J. Devlin, M.-W. Chang, K. Lee, and K. Toutanova, ``BERT: Pre-training of deep bidirectional transformers for language understanding,'' in \textit{Proc. NAACL-HLT}, 2019, pp. 4171--4186.

\bibitem{brooke1996}
J. Brooke, ``SUS: A quick and dirty usability scale,'' in \textit{Usability Evaluation in Industry}, P. W. Jordan \textit{et al.}, Eds. London, U.K.: Taylor \& Francis, 1996, pp. 189--194.

\bibitem{bangor2008}
A. Bangor, P. T. Kortum, and J. T. Miller, ``An empirical evaluation of the System Usability Scale,'' \textit{Int. J. Hum.-Comput. Interact.}, vol. 24, no. 6, pp. 574--594, 2008.

\bibitem{lops2011}
P. Lops, M. de Gemmis, and G. Semeraro, ``Content-based recommender systems: State of the art and trends,'' in \textit{Recommender Systems Handbook}. Boston, MA, USA: Springer, 2011, pp. 73--105.

\bibitem{ricci2015}
F. Ricci, L. Rokach, and B. Shapira, Eds., \textit{Recommender Systems Handbook}, 2nd ed. New York, NY, USA: Springer, 2015.

\bibitem{karpukhin2020}
V. Karpukhin \textit{et al.}, ``Dense passage retrieval for open-domain question answering,'' in \textit{Proc. EMNLP}, 2020, pp. 6769--6781.

\bibitem{cohan2020}
A. Cohan, S. Feldman, I. Beltagy, D. Downey, and D. Weld, ``SPECTER: Document-level representation learning using citation-informed transformers,'' in \textit{Proc. 58th Annu. Meeting Assoc. Comput. Linguistics (ACL)}, 2020, pp. 2270--2282.

\bibitem{singh2023}
A. Singh \textit{et al.}, ``SciRepEval: A multi-format benchmark for scientific document representations,'' in \textit{Proc. Conf. Empirical Methods Natural Lang. Process. (EMNLP)}, 2023.

\bibitem{kinney2023}
R. Kinney \textit{et al.}, ``The Semantic Scholar Open Data Platform,'' \textit{arXiv preprint arXiv:2301.10140}, 2023.

\bibitem{carbonell1998}
J. Carbonell and J. Goldstein, ``The use of MMR, diversity-based reranking for reordering documents and producing summaries,'' in \textit{Proc. 21st Annu. Int. ACM SIGIR Conf. Res. Develop. Inf. Retrieval}, 1998, pp. 335--336.

\bibitem{gemini2023}
Gemini Team, Google, ``Gemini: A family of highly capable multimodal models,'' \textit{arXiv preprint arXiv:2312.11805}, 2023.

\bibitem{sinnott1984}
R. W. Sinnott, ``Virtues of the haversine,'' \textit{Sky and Telescope}, vol. 68, no. 2, p. 159, 1984.
\end{thebibliography}

\end{document}
